\chapter{Code and Math}
\label{chap:code-math}

\section{Code listings}

Code blocks are highlighted by \texttt{minted}, which supports hundreds of
languages. Each listing gets a caption, a number and an entry in the list of
listings. Short inline code works too: \mintinline{python}{len(items) > 0}.

\begin{listing}[H]
  \begin{minted}{python}
from dataclasses import dataclass

@dataclass
class Result:
    name: str
    score: float

def best(results: list[Result]) -> Result:
    """Return the result with the highest score."""
    return max(results, key=lambda r: r.score)
  \end{minted}
  \caption{Python.}
  \label{lst:python}
\end{listing}

\begin{listing}[H]
  \begin{minted}{json}
{
  "name": "my-report",
  "build": "latexmk",
  "outputs": ["main.pdf"]
}
  \end{minted}
  \caption{JSON.}
  \label{lst:json}
\end{listing}

\begin{listing}[H]
  \begin{minted}[linenos=false]{bash}
latexmk -pdf main.tex   # build
latexmk -c              # clean auxiliary files
  \end{minted}
  \caption{Shell commands, without line numbers.}
  \label{lst:bash}
\end{listing}

\section{Equations}

Inline math sits in the text, like $E = mc^2$ or $\sum_{i=1}^{n} x_i$.
Numbered equations stand on their own line and can be referenced, as
\cref{eq:mean,eq:variance}:
\begin{equation}
  \bar{x} = \frac{1}{n} \sum_{i=1}^{n} x_i
  \label{eq:mean}
\end{equation}
\begin{equation}
  \sigma^2 = \frac{1}{n} \sum_{i=1}^{n} \left( x_i - \bar{x} \right)^2
  \label{eq:variance}
\end{equation}

Several equations can be aligned on a symbol with \texttt{align}:
\begin{align}
  f(x) &= (x + 1)^2 \nonumber \\
       &= x^2 + 2x + 1 \label{eq:expanded}
\end{align}

Matrices and cases:
\[
  A = \begin{pmatrix} 1 & 2 \\ 3 & 4 \end{pmatrix},
  \qquad
  |x| = \begin{cases} x & \text{if } x \geq 0 \\ -x & \text{otherwise} \end{cases}
\]

\section{Definitions and theorems}

\begin{definition}[Precision]
  \label{def:precision}
  The precision of a classifier is the share of its positive predictions
  that are correct: $P = \frac{TP}{TP + FP}$.
\end{definition}

\begin{theorem}
  \label{thm:pythagoras}
  In a right triangle with legs $a$ and $b$ and hypotenuse $c$,
  $a^2 + b^2 = c^2$.
\end{theorem}

\begin{proof}
  Arrange four copies of the triangle inside a square of side $a + b$ and
  compare the areas.
\end{proof}

\begin{example}
  With $a = 3$ and $b = 4$, \cref{thm:pythagoras} gives $c = 5$.
\end{example}
